\documentclass[10pt,a4paper]{article}
\usepackage{amsmath,amssymb,amsthm,graphicx}
% Card-local minimal preamble (journal unit sigma-2023-023-spectral-asymmetry-and-index-theory-on-manif, P2-ahl16 precedent).
%% Verbatim from the pinned source: the author theorem environments and macro block,
%% so every macro and environment of the counted spans is available.
% Omitted and why: the documentclass line (the journal class sigma is not on the host;
%% the assembler supplies the class) and page-geometry/font blocks (they fight the
%% forced xelatex engine).
\usepackage{amsmath,amssymb,amsthm,mathtools}


\numberwithin{equation}{section}
\newtheorem{Theorem}{Theorem}[section]
\newtheorem{Corollary}[Theorem]{Corollary}
\newtheorem{Lemma}[Theorem]{Lemma}
\newtheorem{Proposition}[Theorem]{Proposition}
 { \theoremstyle{definition}
\newtheorem{Definition}[Theorem]{Definition}
\newtheorem{Note}[Theorem]{Note}
\newtheorem{Example}[Theorem]{Example}
\newtheorem{Remark}[Theorem]{Remark} }
 \DeclareMathOperator{\Tr}{Tr}
 \DeclareMathOperator{\tr}{tr}
 \DeclareMathOperator{\Str}{Str}
 \DeclareMathOperator{\id}{id}
\DeclareMathOperator{\im}{im}
\DeclareMathOperator{\spec}{spec}
\DeclareMathOperator{\AS}{AS}
\DeclareMathOperator{\HS}{HS}
\DeclareMathOperator{\APS}{APS}
\DeclareMathOperator{\Dir}{Dir}
\DeclareMathOperator{\Neu}{Neu}
\DeclareMathOperator{\Id}{Id}
\DeclareMathOperator{\dom}{dom}
\DeclareMathOperator{\Ei}{Ei}
\DeclareMathOperator{\Idem}{Idem}
 \DeclareMathOperator{\Stab}{Stab}
\DeclareMathOperator{\Ad}{Ad}
\DeclareMathOperator{\ad}{ad}
\DeclareMathOperator{\ind}{index}
\DeclareMathOperator{\ext}{ext}
\DeclareMathOperator{\grad}{grad}
\DeclareMathOperator{\End}{End}
\DeclareMathOperator{\Hom}{Hom}
\DeclareMathOperator{\sgn}{sgn}
\DeclareMathOperator{\ch}{ch}
\DeclareMathOperator{\Todd}{Todd}
\DeclareMathOperator{\topp}{top}
\DeclareMathOperator{\reg}{reg}
\DeclareMathOperator{\mat}{mat}
\DeclareMathOperator{\diag}{diag}
\DeclareMathOperator{\Lie}{Lie}
\DeclareMathOperator{\rank}{rank}
\DeclareMathOperator{\sign}{sign}
\DeclareMathOperator{\diam}{diam}
\DeclareMathOperator{\red}{red}
\DeclareMathOperator{\LIM}{LIM}
\DeclareMathOperator{\Imag}{Im}
\DeclareMathOperator{\Real}{Re}
\DeclareMathOperator{\Spin}{Spin}
\DeclareMathOperator{\OO}{O}
\DeclareMathOperator{\SO}{SO}
\DeclareMathOperator{\SL}{SL}
\DeclareMathOperator{\GL}{GL}
\DeclareMathOperator{\SU}{SU}
\DeclareMathOperator{\U}{U}
\DeclareMathOperator{\Cl}{Cl}
 \DeclareMathOperator{\Int}{int}
 \DeclareMathOperator{\Ind}{Ind}
 \DeclareMathOperator{\ev}{ev}
 \DeclareMathOperator{\Res}{Res}
 \DeclareMathOperator{\Thom}{Thom}
 \DeclareMathOperator{\Span}{span}
 \DeclareMathOperator{\supp}{supp}
\DeclareMathOperator{\pt}{pt}
 \DeclareMathOperator{\even}{even}
 \DeclareMathOperator{\odd}{odd}
 \DeclareMathOperator{\DInd}{D-Ind}
\DeclareMathOperator{\KInd}{K-Ind}
\DeclareMathOperator{\vol}{vol}
\newcommand{\Spinc}{\Spin^c}
 \newcommand{\R}{\mathbb{R}}
 \newcommand{\C}{\mathbb{C}}
 \newcommand{\N}{\mathbb{N}}
 \newcommand{\Z}{\mathbb{Z}}
 \newcommand{\Q}{\mathbb{Q}}
 %\newcommand{\K}{\mathbb{K}}
 \newcommand{\bT}{\mathbb{T}}
 \newcommand{\bP}{\mathbb{P}}
\newcommand{\g}{\mathfrak{g}}
\newcommand{\h}{\mathfrak{h}}
\newcommand{\p}{\mathfrak{p}}
\newcommand{\kg}{\mathfrak{g}}
\newcommand{\kt}{\mathfrak{t}}
\newcommand{\ka}{\mathfrak{a}}
\newcommand{\XX}{\mathfrak{X}}
\newcommand{\kh}{\mathfrak{h}}
\newcommand{\kp}{\mathfrak{p}}
\newcommand{\kk}{\mathfrak{k}}
\newcommand{\kn}{\mathfrak{n}}
\newcommand{\cA}{\mathcal{A}}
\newcommand{\cE}{\mathcal{E}}
\newcommand{\calL}{\mathcal{L}}
\newcommand{\calH}{\mathcal{H}}
\newcommand{\cO}{\mathcal{O}}
\newcommand{\cB}{\mathcal{B}}
\newcommand{\cK}{\mathcal{K}}
\newcommand{\cP}{\mathcal{P}}
\newcommand{\cN}{\mathcal{N}}
\newcommand{\calD}{\mathcal{D}}
\newcommand{\cC}{\mathcal{C}}
\newcommand{\calS}{\mathcal{S}}
\newcommand{\cM}{\mathcal{M}}
\newcommand{\cU}{\mathcal{U}}
\newcommand{\cF}{\mathcal{F}}
\newcommand{\cCM}{\cC}
\newcommand{\PM}{P}
\newcommand{\DM}{D}
\newcommand{\LM}{L}
\newcommand{\vM}{v}
\newcommand{\sumGam}{\textstyle{\sum_{\Gamma}} }
\newcommand{\sigDg}{\sigma^D_g}
\newcommand{\Bigwedge}{\textstyle{\bigwedge}}
\newcommand{\ii}{\sqrt{-1}}
\newcommand{\Ubar}{\overline{U}}
\usepackage{cleveref}
\crefname{originalequation}{original equation}{original equations}
\Crefname{originalequation}{Original equation}{Original equations}
\makeatletter
\newlabel{eq D0+ P}{{3.20}{source}{}{sourcepointer-1}{}}
\newlabel{eq D0+ P@cref}{{[originalequation][0][]3.20}{[1][0][]source}{}{}{}}
\makeatother
\crefname{originalequation}{original equation}{original equations}
\Crefname{originalequation}{Original equation}{Original equations}
\makeatletter
\newlabel{eq D0- 1-P}{{3.21}{source}{}{sourcepointer-2}{}}
\newlabel{eq D0- 1-P@cref}{{[originalequation][0][]3.21}{[1][0][]source}{}{}{}}
\makeatother
\title{The two first-order model operators are each other's adjoints when h has a positive lower bound}
\author{Hochs, Peter, Wang, Hang}
\date{doi:10.3842/SIGMA.2023.023; CC BY 4.0}
\begin{document}
\maketitle
\setcounter{section}{3}
\setcounter{equation}{32}
\setcounter{Theorem}{13}
\noindent\textit{Editorial note.} All delivered mathematical text is the author's own, copied from the pinned journal source file sources/sigma-2023-023-spectral-asymmetry-and-index-theory-on-manif/pinned-original-source.tex (Spectral Asymmetry and Index Theory on Manifolds with Generalised Hyperbolic Cusps); no statement or proof step was rewritten, reordered or removed.\par\medskip
\noindent\textit{Editorial note.} The counted proof is the author's complete argument: the counted statement and proof are the author's verbatim text; the external references eq D0+ P and eq D0- 1-P are navigation stubs carrying the article's printed numbers; the display labels eq adjoints DC 1 and eq adjoints DC 2 lie inside the span and are internal to the card.\par\medskip
\noindent\textit{Editorial note.} The label prop adjoints DC names the counted statement, and the proof environment immediately following it in the source is delivered verbatim as its proof.\par\medskip
\noindent\textit{Editorial note.} The counted statement is carried by the author's own Proposition environment in the article source.\par\medskip
\noindent\textit{Editorial note.} Journal-unit assembly: card-local preamble (author theorem environments and macros verbatim), article class, no bibliography in the counted span (the span carries no \textbackslash{}cite commands); references leaving the span are declared as navigation stubs carrying the article's printed numbers.\par\medskip
\noindent\textit{Editorial note.} COUNTER REPAIR 2026-10-14: the r001 seeds were derived from the census v3.5+ printed field, which counts \textbackslash{}section and \textbackslash{}subsection jointly (subsection-inclusive) and diverges from the article's own numbering; the seeds above are recomputed by the lane's true-numbering rule (\{"section": 3, "equation": 32, "Theorem": 13\}).\par\medskip
\noindent\textit{Editorial scope.} the counted statement and its own proof environment only\par\medskip
\noindent\textit{Allowed standard prerequisites (editorial).}\begin{itemize}
\item The two first-order model operators and their domains. The operators eq D0+ P and eq D0- 1-P acting on L\textasciicircum{}2(tilde S\_C\textasciicircum{}+), with the positive function h on the cylinder.. The pair the proposition relates; declared as the article's standing notation, no external stub beyond the two operator-display stubs.
\end{itemize}\medskip
\ifdefined\hypertarget\hypertarget{sourcepointer-1}{}\fi
\noindent\textit{Source reference.} Original equation 3.20: sources/sigma-2023-023-spectral-asymmetry-and-index-theory-on-manif/pinned-original-source.tex, source lines 1005--1005. Navigation pointer invoked at line(s) 1199: the counted text cites that numbered equation of the same article, which is not retained in this card, so the reference is declared as a pointer carrying the number the article itself prints for it (the section and the place of the equation inside that section, derived by the census v3.1 rule for the section-scoped equation counter); the cited equation is an allowed input and none of it is reproduced inside the span.\par\medskip
\ifdefined\hypertarget\hypertarget{sourcepointer-2}{}\fi
\noindent\textit{Source reference.} Original equation 3.21: sources/sigma-2023-023-spectral-asymmetry-and-index-theory-on-manif/pinned-original-source.tex, source lines 1006--1006. Navigation pointer invoked at line(s) 1199: the counted text cites that numbered equation of the same article, which is not retained in this card, so the reference is declared as a pointer carrying the number the article itself prints for it (the section and the place of the equation inside that section, derived by the census v3.1 rule for the section-scoped equation counter); the cited equation is an allowed input and none of it is reproduced inside the span.\par\medskip
\begin{Proposition}\label{prop adjoints DC}
If $h$ has a positive lower bound, then
the two operators~\eqref{eq D0+ P} and~\eqref{eq D0- 1-P} are each other's adjoints.
\end{Proposition}

\begin{proof}
We claim that for all $s_P \in W^1_{\widetilde{D}_C}\big(\widetilde{S}_C^+; P\big)$ and $s_{1-P} \in W^1_{\widetilde{D}_C}\big(\widetilde{S}_C^+; 1-P\big)$,
\begin{equation}
\label{eq adjoints DC 1}
\big(\widetilde{D}_C^+ s_P, s_{1-P}\big)_{L^2} = \big(s_P, \widetilde{D}_C^- s_{1-P}\big)_{L^2}.
\end{equation}
Indeed, suppose that
\[
s_P = s_P^N \otimes f_{P},\qquad s_{1-P} = s_{1-P}^N \otimes f_{1-P},
\]
for $s_P^N, s_{1-P}^N \in L^2\big(S^+|_N\big)$ and $f_P, f_{1-P} \in L^2[0, \infty)$ such that $s_P \in W^1_{\widetilde{D}_C}\big(\widetilde{S}_C^+; P\big)$ and $s_{1-P} \in W^1_{\widetilde{D}_C}\big(\widetilde{S}_C^+; 1-P\big)$. Then by self-adjointness of $D_N^+$ and integration by parts, the left-hand side of~\eqref{eq adjoints DC 1} equals
\begin{equation} \label{eq adjoints DC 2}
\big(s_P^N, D_N^+ s_{1-P}^N\big)_{L^2(S^+|_N)} \bigg(
\big(f_P \overline{f_{1-P}}\big)|_{0}^{\infty} + \int_0^{\infty} f_P(x)\bigl(-\overline{f_{1-P}'(x)} +h(x) \overline{f_{1-P}(x)}\bigr)\, {\rm d}x\bigg).
\end{equation}
If we further decompose the expressions with respect to eigenspaces of $D_N^+$, and use that the components of $f_P$ for positive eigenvalues equal zero at zero, and the components of $f_{1-P}$ for negative eigenvalues equal zero at zero, then we find that the components for all eigenvalues of $f_P \overline{f_{1-P}}$ are zero at zero. So the term $\big(f_P \overline{f_{1-P}}\big)|_{0}^{\infty} $ in~\eqref{eq adjoints DC 2} vanishes, and~\eqref{eq adjoints DC 2} equals the right-hand side of~\eqref{eq adjoints DC 1}.

Now if $\sigma_P := \widetilde{D}_C^+ s_P$ and $\sigma_{1-P} := \widetilde{D}_C^- s_{1-P}$, then
\eqref{eq adjoints DC 1} becomes
\[
\big( \sigma_P, \big(\widetilde{D}_C^-\big)^{-1} \sigma_{1-P}\big)_{L^2} = \big( \big(\widetilde{D}_C^+\big)^{-1}\sigma_P, \sigma_{1-P}\big)_{L^2}.
\]
The sections $\sigma_P$ and $\sigma_{1-P}$ of this type are dense in $L^2\big(\widetilde{S}_C^+\big)$. And the inverse operators $\big(\widetilde{D}_C^+\big)^{-1}$ and $\big(\widetilde{D}_C^-\big)^{-1}$ are bounded, so we find that
\[
\big(\big(\widetilde{D}_C^+\big)^{-1}\big)^* = \big(\widetilde{D}_C^-\big)^{-1}.
\]
This implies that $\big(\widetilde{D}_C^+\big)^* = \widetilde{D}_C^-$.
\end{proof}

\end{document}
